\subsection{记号}

给定一个实函数 \(f \in \mathcal{H}(\mathfrak{F}, \mathbf{R})\) ，在一个公式中用 \(O(f)\) 表示一个受 \(f\) 控制的函数、用 \(o(f)\) 表示一个相对于 \(f\) 可忽略的函数，这往往是方便的。当在一个证明中出现多个受同一个函数 \(f\) 控制的（相应地，相对于 \(f\) 可忽略的）函数时，我们把它们记作 \(O_1(f)\) 、 \(O_2(f)\) 等（相应地， \(o_1(f)\) 、 \(o_2(f)\) 等）。

许多作者不加区分地把一个证明中所有受 f 控制的（相应地，相对于 f 可忽略的）函数都记作 \(O(f)\) （相应地， \(o(f)\) ），这是一种可能引起混淆的语言滥用。

用这个记号，我们可以把命题 1、2、3 (V, p.~213) 表述如下：若 \(g = O_1(f)\) 且 \(h = O(g)\) ，则 \(h = O_2(f)\) ；我们可以写

\[
\sum_ {i = 1} ^ {n} \lambda_ {i} O _ {i} (f) = O _ {n + 1} (f) \quad (\lambda_ {i} \text {   scalars })\tag{1}
\]

\[
O (f) O (g) = O (f g).\tag{2}
\]

类似地，命题 4 (V, p.~215) 表明，若 \(g = O_1(f)\) 且 \(h = o(g)\) （相应地， \(g = o_1(f)\) 且 \(h = O(g)\) ），则 \(h = o_2(f)\) ，而命题 5 与 6 (V, p.~5) 可以表述成形式

\[
\sum_ {i = 1} ^ {n} \lambda_ {i} o _ {i} (f) = o _ {n + 1} (f) \quad (\lambda_ {i} \text {   scalars })\tag{3}
\]

\[
o (f) O (g) = o (f g).\tag{4}
\]

关系 \(f \sim g\) 等价于 \(f = g + o(g)\) 。记号 \(O(1)\) （相应地， \(o(1)\) ）表示在 \(\mathfrak{F}\) 的某个集上有界的一个函数（相应地，沿 \(\mathfrak{F}\) 趋于 0 的一个函数）。

